%% kapa: exam=1 round=36 year=2025 period=2 form=A subject=appraisal_law subject_ko=감정평가관계법규 seq=20 number=20
%% source: Q-Net 감정평가사 기출문제 (cst00309) artlSeq=5231525 file=raw/1008/5231525/2025년 제36회 감정평가사 1차시험 2교시 문제지 원본.pdf
%% conversion: kordoc 4.24.51 -> tools/gen1.py (auto); review=proofread-sonnet; flags=none
\begin{question}{20}
국유재산법령상 행정재산에 관한 설명으로 옳지 않은 것은?
\begin{choices}
  \item 행정재산의 종류에는 공용재산, 공공용재산, 기업용재산, 보존용재산이 있다.
  \item 총괄청은 기부채납에 따른 행정재산의 취득에 관한 사무를 중앙관서의 장에게 위임한다.
  \item 행정재산은 ｢민법｣에도 불구하고 시효취득의 대상이 되지 아니한다.
  \item 총괄청은 사용을 승인한 행정재산의 재산관리가 감사 결과 부당한 것으로 인정된 경우 국유재산정책심의위원회의 심의를 거쳐 그 사용 승인을 철회할 수 있다.
  \item 중앙관서의 장은 행정재산으로 기부하는 재산에 대하여 기부자에게 무상으로 사용허가하여 줄 것을 조건으로 그 재산을 기부하는 경우 그 기부를 받아서는 아니 된다.
\end{choices}
\end{question}
\answer{5}
